% Standalone manuscript: upload this complete file as main.tex to Overleaf.
% Bibliography is embedded; no external figures, .bib files, or custom class.
\documentclass[11pt,a4paper]{article}
\usepackage[T1]{fontenc}
\usepackage[utf8]{inputenc}
\usepackage{lmodern}
\usepackage[margin=24mm]{geometry}
\usepackage{amsmath,amssymb,amsfonts}
\usepackage{array,booktabs}
\usepackage{microtype}
\usepackage{xurl}
\usepackage[hidelinks]{hyperref}
\setlength{\emergencystretch}{3em}
\setlength{\parskip}{0.2em}
\setlength{\parindent}{1.2em}
\setlength{\abovedisplayskip}{7pt plus 2pt minus 2pt}
\setlength{\belowdisplayskip}{7pt plus 2pt minus 2pt}
\renewcommand{\arraystretch}{1.08}
\newcommand{\R}{\mathbb R}
\newcommand{\indic}{\mathbb I}
\newcommand{\clip}{\operatorname{clip}}
\title{Two-View Image-Based Anthropometry:\\
Profile-Guided Reconstruction and Numerical Verification}
\author{Internal review draft}
\date{}
\begin{document}
\maketitle

\begin{abstract}
This paper specifies a two-view anthropometry method and examines its numerical components with reproducible synthetic tests. Pretrained perception supplies frontal and profile silhouettes calibrated by stature. Filtered profile intervals are coupled to frontal column thickness to constrain depth at narrow structures; a bounded waist rule applies local occupancy deletion. Measurements use explicitly identified reconstruction branches and defined section levels. A centred, scaled least-squares reference estimator recovers elliptical sections and distinguishes unavailable sensitivity from a computed zero. Across 27 noiseless ellipse cases spanning rotations, translations, and coordinate scales, its maximum relative semi-axis error was $3.35\times10^{-11}$. Controlled geometric experiments examine the profile constraint and its failure cases. These results verify specified computations under synthetic assumptions. Independent anthropometric agreement, the anatomical value of the waist rule, and improvement over established two-view methods remain open empirical questions.
\end{abstract}
\noindent\textbf{Keywords:} image-based anthropometry; visual hull; profile-guided reconstruction; geometric measurement.

\section{Introduction and Related Work}
Photographic anthropometry requires correspondence between image-derived geometry and a specified physical measurement. ISO~8559-1 distinguishes circumferences, breadths, and surface paths used in clothing measurement~\cite{iso8559}. Meunier and Yin evaluated a two-dimensional photographic measurement system~\cite{meunier2000}; Hung \emph{et al.} estimated circumferences from image-derived widths and depths~\cite{hung2004}. Automatic measurement from frontal and lateral silhouettes was subsequently developed by Aslam \emph{et al.}~\cite{aslam2017}. These studies establish substantial prior work for two-view measurement and elliptical section models.

The visual hull represents the volume consistent with observed silhouettes~\cite{laurentini1994visual}. Human-shape methods also include the statistical SMPL representation~\cite{loper2015smpl}, volumetric inference in BodyNet~\cite{varol2018bodynet}, and learned reconstruction from orthographic silhouettes~\cite{ji2019}. The present pipeline uses pretrained perception followed by geometric rules with explicit physical units. Its principal proposed adaptation couples a filtered profile interval to frontal run topology and a columnwise thickness bound. This construction addresses depth replication when a profile silhouette contains arms or displaced limbs. A spatially tapered waist correction provides a further local shape constraint. The shoulder model specializes an established ellipse-perimeter approximation to a defined surface path. These adaptations require ablation to establish an accuracy benefit; the supporting segmentation, fitting, and meshing algorithms retain their established attribution.

At the operation level, Hung \emph{et al.} combine photographic breadth and depth with ellipse-based circumference estimation~\cite{hung2004}; Aslam \emph{et al.} identify silhouette fiducials and use an ellipsoid model under a calibrated capture arrangement~\cite{aslam2017}. The present construction instead specifies the occupancy interval assigned to each frontal column before section fitting. Ji \emph{et al.} learn frontal/lateral silhouette features to recover a body-shape representation~\cite{ji2019}. More recent compact multi-view work combines RGB and depth-assisted imagery with SMPL estimation and exemplar optimization~\cite{chen2025compact}. These different reconstruction objectives preclude an accuracy comparison from the present synthetic tests. The proposed contribution is the explicit profile/run-thickness coupling and its measured geometric consequences in the tested cases. The waist deletion is an additional heuristic to evaluate. This positioning does not establish priority for all combinations of two-view silhouettes and shape priors.

The formulation below separates the historical pipeline operations from a corrected conic research reference supplied with this manuscript. The reference is an isolated numerical implementation; it has not been integrated into an end-to-end human evaluation. Chest-level girth uses the profile-guided branch, whereas waist-proxy and hip-landmark girths use a waist-adjusted conventional hull. A separately rounded profile-guided volume supplies the display mesh. Each operation imposes a different section model. Geometric processing uses fixed rules after pretrained perception.

\section{Image Preparation and Metric Alignment}
\subsection{Acquisition, masks, and landmarks}
Inputs are frontal and approximately orthogonal profile RGB images, together with stature $H$ in centimetres. The full crown and feet must be visible; close-fitting clothing and consistent posture reduce silhouette ambiguity. Let $X$ denote lateral position, $Y$ distance downwards from the crown, and $Z$ sagittal position, all in millimetres. A pinhole camera gives $u=f_xx_c/z_c+c_u$ and $v=f_yy_c/z_c+c_v$. Replacing depth by a representative $z_0$ yields the weak-perspective approximation used below. A depth displacement $\Delta z$ changes magnification by approximately $-\Delta z/z_0$. The implementation estimates neither camera intrinsics nor residual yaw, so stature calibration alone does not eliminate perspective error.

For planned acquisition, the camera should remain level with the subject centre, retain its focal setting and subject distance between views, and record distance, image orientation and intended $90^\circ$ view separation. A frontal pose with approximately horizontal, separated arms is the intended domain of the column-thickness rule. Crown/foot contact with an image border, foreground lost after alignment, missing required landmarks and substantial discarded components require recorded rejection or review. These are proposed capture checks; no validated angular or distance tolerance is claimed. For a point satisfying $|\Delta z|\leq d<z_0$, the pinhole magnification relative to the reference plane differs by at most $d/(z_0-d)$. This conditional geometric bound quantifies one perspective effect and does not bound the error of the complete anthropometry method.

The implementation uses MediaPipe~\cite{lugaresi2019mediapipe} and its Pose Landmarker segmentation output~\cite{mediapipeguide}. The heavy pretrained model operates in single-image, single-person mode with detection and presence thresholds of $0.5$. Contiguous RGB arrays are padded on the right with replicated columns to a width divisible by four; the segmentation map is cropped back afterwards. This alignment step preserves image scale. A mismatched map is resized bilinearly. Foreground is defined by the strict probability threshold $p>0.5$, followed by selection of the largest eight-connected component. Enclosed background holes smaller than $2\%$ of foreground area are filled; border-connected background is retained. Missing pose or segmentation output, or foreground occupying less than $2\%$ of the frame, causes rejection. Default mask dilation is zero.

An optional GrabCut baseline uses the foreground/background colour-mixture and boundary formulation of Rother \emph{et al.}~\cite{rother2004grabcut}. Landmark-based initialization constructs a torso polygon, limb capsules, and a head region. With estimated pixel stature $H_p=(y_{\rm ankle}-y_{\rm shoulder})/0.779$, seed erosion is $0.012H_p$ and the probable-foreground and background dilations are $0.045H_p$ and $0.13H_p$. A border ring spans $2\%$ of the smaller image dimension. Five iterations with random seed zero precede a $7\times7$ elliptical closing and the same component cleanup. This baseline is separate from the default learned mask.

Landmarks are refined by a second inference on the silhouette bounding box expanded by $15\%$. Crops whose longer dimension is below 1024 pixels are enlarged bicubically; larger crops retain their resolution. For resize factor $\gamma$, resized crop dimensions $(W_r,H_r)$, and crop origin $(x_0,y_0)$, a normalized landmark $(\xi,\eta)$ maps to $(x_0+W_r\xi/\gamma,\,y_0+H_r\eta/\gamma)$. Shoulder, hip, and ankle pairs (indices 11--12, 23--24, 27--28) determine alignment. Elbows and wrists provide subsequent overlap checks. The fit is rejected when more than three requested landmarks have visibility below $0.3$.

\subsection{Weighted alignment and calibration}
For image points $p_i$, canonical targets $q_i$ place each shoulder, hip, and ankle pair horizontally while preserving its Euclidean pair distance and mean vertical position. Shoulder and hip midpoints define a common spine coordinate; ankle stance and image ordering are preserved. Visibility weights are clipped to $[0.05,1]$ and normalized to sum to one. The similarity fit extends orthogonal Procrustes alignment~\cite{schonemann1966generalized}:
\begin{equation}
 \min_{\sigma>0,R\in SO(2),t}\sum_iw_i\|\sigma Rp_i+t-q_i\|^2.
 \label{eq:procrustes}
\end{equation}
Writing $\bar p=\sum_iw_ip_i$, $\bar q=\sum_iw_iq_i$, and
$K=\sum_iw_i(p_i-\bar p)(q_i-\bar q)^T=U\Sigma V^T$, the solution is
\begin{equation}
 \begin{aligned}
 S&=\operatorname{diag}(1,\det(VU^T)),& R&=VSU^T,\\
 \sigma&=\frac{\operatorname{tr}(\Sigma S)}{\sum_iw_i\|p_i-\bar p\|^2},&
 t&=\bar q-\sigma R\bar p.
 \end{aligned}
\end{equation}
Eliminating translation centres the objective; maximizing its cross term gives $R$, and differentiation with respect to scale gives $\sigma$. Near-zero source spread uses unit scale. Rotation beyond $35^\circ$ or scale outside $[0.6,1.6]$ causes rejection. Masks are warped on the original canvas with bilinear interpolation and rethresholded at $0.5$; landmarks receive the same transform. This global alignment can leave articulation differences and can clip foreground at the canvas boundary.

Calibration follows warping. In each view $v\in\{f,s\}$, let $t_v,b_v$ be the first and last occupied rows and $L_v=b_v-t_v+1$. Set $H_m=10H$ and $s_v=H_m/L_v$ millimetres per pixel. The horizontal origin $c_v$ is the equally weighted mean of nonempty-row silhouette midpoints. Projection is
\begin{equation}
 (u_f,v_f)=(c_f+X/s_f,\,t_f+Y/s_f),\qquad
 (u_s,v_s)=(c_s+Z/s_s,\,t_s+Y/s_s).
 \label{eq:projection}
\end{equation}
Nearest-pixel mask lookup returns zero outside the image. Occupancy arrays use $(X,Y,Z)$ ordering. OpenCV section coordinates are $(\text{column},\text{row})=(Z\text{-index},X\text{-index})$ and must be reversed when recovering metric $(X,Z)$ points. Supplied stature is a common scale input to both views. Continuous geometric measurements therefore satisfy $\Delta C/C\simeq\Delta H/H$; fixed metric thresholds and voxel discretization can change the discrete response.

\section{Reconstruction and Local Shape Constraints}
\subsection{Conventional occupancy and branch selection}
For binary masks $M_f,M_s$, the conventional orthographic occupancy is
\begin{equation}
 O_0(X,Y,Z)=M_f(u_f,v_f)M_s(u_s,v_s).
 \label{eq:hull}
\end{equation}
At a given height, every frontal run pairs with every profile run. With $n$ frontal and $m$ profile intervals, this produces up to $nm$ rectangular regions, including limb-depth replication. Reconstruction first uses a 10\,mm grid over silhouette extrema with 30\,mm lateral and sagittal margins and $0\leq Y<H_m$. The occupied coarse bounding box is expanded by one coarse cell and resampled uniformly at $h=2$\,mm. Features absent from the coarse sample can consequently be missed. The fine stage stores a dense array over this box.

The largest 26-connected volume component is retained. Within successive $XZ$ slices, eight-connected components are tracked greedily by nearest centroid, with a maximum linking distance of 60\,mm. Tracks spanning at least five slices have their centroid coordinates and four half-extents filtered by a cubic Savitzky--Golay fit~\cite{savitzky1964smoothing}, using at most 15 slices. Negative half-extents are clipped to zero. Each original component is intersected with its smoothed bounding rectangle, and the largest volume component is selected again. This operation trims occupancy; it cannot recover missing material.

The resulting conventional hull receives the waist correction below and supplies waist and hip sections. Chest uses an independently constructed, unrounded profile-guided volume, followed by its own section conditioning. A rounded copy of the profile-guided volume is used for display. Failure to construct a usable profile core sends chest and display processing to the corrected conventional hull. Shoulder and hip-to-ankle outputs use masks and landmarks directly.

\subsection{Profile filtering and columnwise depth}
At profile row $y$, select the widest occupied run $[\ell_y,u_y]$, with width $d_y=u_y-\ell_y$ and centre $c_y=(\ell_y+u_y)/2$. Reject missing runs and rows satisfying
\begin{equation}
 d_y>1.8\operatorname{median}(d)\quad\text{or}\quad
 |c_y-\operatorname{median}(c)|>
 \max\{6\operatorname{MAD}(c),40\text{ pixels}\}.
\end{equation}
Here MAD is the unscaled median absolute deviation. A 51-row vertical dilation expands the rejected-row mask by approximately 25 rows in either direction. Accepted endpoints are linearly interpolated, with endpoint clamping, then filtered by a 101-row median and a cubic Savitzky--Golay window of at most 61 odd-numbered rows. Windows shrink or are bypassed for short sequences. No accepted rows causes branch failure. For offsets $r$ in a filtering window, the local polynomial coefficients solve
\begin{equation}
 \hat a=(V^TV)^{-1}V^Tf,\qquad V_{rk}=r^k,\quad k=0,\ldots,3;
 \qquad \hat f(0)=\hat a_0.
\end{equation}
Window lengths are in pixels, so their physical bandwidth depends on $s_s$. Interpolated heights are flagged because rejected observations have been replaced by a smooth depth prior.

Converting filtered endpoints to metric coordinates gives $\zeta_-(Y)$ and $\zeta_+(Y)$. For each fixed frontal column $X$, identify its contiguous vertical foreground runs. If the run containing $Y$ spans $n_Y(X,Y)$ fine-grid samples, define
\begin{equation}
 \begin{aligned}
 T(X,Y)&=h\,n_Y(X,Y),&
 D_p(Y)&=\zeta_+(Y)-\zeta_-(Y),\\
 z_c(Y)&=\tfrac12[\zeta_+(Y)+\zeta_-(Y)],&
 z_\pm(X,Y)&=z_c(Y)\pm\tfrac12\min\{D_p(Y),T(X,Y)\},\\
 O_p(X,Y,Z)&=M_f(u_f,v_f)\indic\{z_-\leq Z\leq z_+\}.&&
 \end{aligned}
 \label{eq:profile}
\end{equation}
The profile supplies a shared sagittal midpoint and available depth. The frontal vertical run supplies a local thickness bound: an approximately horizontal limb has a short vertical extent at its outer columns, even when it joins the torso horizontally. This coupling is the main proposed reconstruction adaptation. Its thickness and midpoint assumptions can misrepresent sloping limbs or displaced anatomy. Endpoint interpolation may also extend occupancy beyond the observed profile mask; $O_p$ is therefore a model-conditioned reconstruction. Computation proceeds in chunks of 64 frontal columns while retaining the dense output volume.

For a straight strip of perpendicular thickness $t$ at angle $\varphi$ to the horizontal, an interior vertical intersection has length $t/|\cos\varphi|$ before finite-length effects. Thus the cap approximates thickness only near a horizontal orientation; its restriction weakens as the limb becomes vertical. A shared midpoint cannot recover a limb's unknown displacement in $Z$. The intended use therefore requires near-horizontal arms and reports interpolation at each measurement level. An oblique pose, a displaced limb or an interpolated torso band remains a model-assumption case even when reconstruction completes.

\subsection{Measurement levels and curvature-dependent waist correction}
The frontal torso trace starts from the widest horizontal run and subsequently selects the run nearest the preceding midpoint. Separate left/right arm-separation tests use five-row margins. Missing trace values are interpolated and smoothed by a cubic 21-row fit. With mean shoulder and hip levels $Y_s,Y_h$, the waist proxy is $Y_w=Y_s+0.688(Y_h-Y_s)$, snapped to the nearest available trace row. The factor approximates $(0.818-0.620)/(0.818-0.530)$ from engineering body-height fractions. Chest is evaluated at $Y_c=Y_s+0.35(Y_w-Y_s)$; hip uses the landmark level. These prescribed levels do not search for maximum chest or hip girth.

Each torso half-width is clipped using its 15th percentile in a $\pm40$\,mm neighbourhood; nonpositive widths are rejected. Girth extraction expands these bounds by a factor of $1.15$, whereas shoulder extraction uses the unexpanded trace. Conventional-hull sections also receive a centred depth cap $D_c=\min(D,0.75W)$ when a profile elbow or wrist lies within 90\,mm of the measurement level, or when those landmarks are unavailable. Removal exceeding $1\%$ of original depth generates an advisory. Profile-guided chest receives only lateral torso clipping.

Waist refinement locates a local minimum of each half-width within $\pm15$ trace samples of $Y_w$, then walks along rising widths to obtain an enclosing interval $[Y_0,Y_1]$. Fit $w(Y)\simeq\alpha(Y-Y_w)^2+\beta(Y-Y_w)+\gamma$. The graph curvature is $2\alpha/(1+\beta^2)^{3/2}$; the implementation uses the small-slope magnitude $\kappa=|2\alpha|$. Writing $L=Y_1-Y_0$, the inward displacement is
\begin{equation}
 \begin{aligned}
 R_c&=\kappa^{-1},\qquad
 \theta=2\arcsin\!\left(\clip\!\left(\frac{L}{2R_c},-0.999,0.999\right)\right),\\
 d&=\min\!\left\{\frac{R_c}{2}[1-\cos(\theta/2)],\,0.40w_{\rm loc}\right\}.
 \end{aligned}
 \label{eq:waist}
\end{equation}
Here $w_{\rm loc}$ is the smaller of the endpoint and local waist half-widths. The correction is skipped for prominence below 2\,mm, nonpositive $L$, or $\kappa<10^{-6}$\,mm$^{-1}$. The displacement uses half the circular sagitta; before clipping, its small-curvature limit is $d\simeq\kappa L^2/16$. This attenuated geometric prior is derived from the frontal width trace.

Within $Y_w\pm15$\,mm, sampled Hann weights $\omega_j=\tfrac12[1-\cos(2\pi j/(n+1))]$, $j=1,\ldots,n$, suppress abrupt vertical changes. Sagittal tapering uses $t(z)=\cos[\clip(\pi z/(2z_h),-\pi/2,\pi/2)]$, where $z_h$ is the maximum occupied $|Z|$, bounded below by 1\,mm. At each $(Y,Z)$, the left and right torso limits move inward by $\omega_jd_Lt(z)$ and $\omega_jd_Rt(z)$. Only occupancy inside the traced torso interval expanded by $1.2$ is eligible for deletion; exterior material is retained. Since $0\leq\omega_jt(z)\leq1$ and $0\leq d\leq0.40w_{\rm loc}$, the displacement is bounded and the adjusted occupancy is a subset of the input. A zero displacement is a no-op and cannot add section area. Actual section-area change is $\Delta A=h^2(N_{\rm before}-N_{\rm after})\geq0$. These bounds establish the operation's direction, not the correctness of the deleted anatomy. The attenuation factor, cap and 15\,mm band are design parameters; global $Z=0$ is the taper origin and can give an asymmetric effect on an off-centre section. The rule is a bounded occupancy heuristic, with no anatomical identification theorem.

\section{Section Measurement and Surface Export}
\subsection{Conditioning and conic recovery}
The nearest $XZ$ slice is selected for each measurement level, giving a level quantization error no greater than $h/2$ within the sampled range. After torso clipping and any conditional depth cap, an OpenCV elliptical closing with nominal radius 4\,mm fills small gaps. The largest eight-connected section component is retained and opened with radius $r=0.9\min(W/2,D/2)$, rounded to grid units with a minimum of two pixels. In set notation, closing is $(S\oplus B)\ominus B$ and opening is $(S\ominus B)\oplus B$. The latter imposes substantial section rounding, which influences the reported circumference.

All contour points are retained, and the contour of largest area is converted to metric $(X,Z)$ coordinates. A Graham scan~\cite{graham1972efficient} deduplicates points, sorts by polar angle, retains the farther collinear point, and removes non-left turns. Complexity is $O(n\log n)$. At least five hull vertices are required. Convexification discards section concavities before fitting.

Let $p_i$ be hull vertices, $ar p$ their mean, and
$\rho=[n^{-1}\sum_i\|p_i-\bar p\|^2]^{1/2}$ their isotropic coordinate scale.
The corrected reference uses $v_i=(p_i-\bar p)/\rho=(\xi_i,\eta_i)^T$ and solves
\begin{equation}
 \min_\beta\|M\beta+\mathbf1\|^2,\qquad
 M_i=(\xi_i^2,\xi_i\eta_i,\eta_i^2,\xi_i,\eta_i).
 \label{eq:fit}
\end{equation}
An SVD-backed least-squares solve avoids forming normal equations. The design must have rank five at the solver's relative rank threshold. Nonfinite coordinates, fewer than five distinct points and zero coordinate spread are rejected. If $\beta=(\beta_1,\ldots,\beta_5)^T$, physical coefficients are
$\phi=(\beta_1/\rho^2,\beta_2/\rho^2,\beta_3/\rho^2,\beta_4/\rho,\beta_5/\rho)^T=(A,B,C,D,E)^T$.
They describe the fixed-constant conic $Ax^2+Bxz+Cz^2+Dx+Ez+1=0$ in millimetres relative to $\bar p$. Equal hull-vertex weights are used. An unrestricted algebraic fit may be a non-ellipse; ellipse-specific constrained fitting provides a relevant alternative~\cite{fitzgibbon1999direct}.
\begin{equation}
 Q=\begin{pmatrix}A&B/2\\B/2&C\end{pmatrix},\quad l=\binom DE,
 \quad \mu=-\tfrac12Q^{-1}l,\quad
 F=\tfrac14l^TQ^{-1}l-1.
 \label{eq:recover}
\end{equation}
Completing the square gives $(u-\mu)^TQ(u-\mu)=F$. The reference solves $Qq=l$ rather than explicitly forming $Q^{-1}$. Eigenvectors determine orientation; the positive ratios $F/\lambda_i$ give squared semi-axes, sorted $a\geq b>0$, with centre $\bar p+\mu$. The quadratic form must be definite, and each eigenvalue magnitude must exceed $10^{-12}\|Q\|_2$. Invalid axis ratios cause rejection. The scaled design condition number is recorded; acceptance of a fit supplies no accuracy guarantee at untested conditioning or incomplete arcs. Residual RMS is dimensionless. The fitted ellipse can extend beyond its source section.

Ramanujan's approximation~\cite{ramanujan1914modular} supplies the perimeter in millimetres,
\begin{equation}
 P(a,b)=\pi[3(a+b)-\sqrt{(3a+b)(a+3b)}],\qquad C=P(a,b)/10\text{ cm}.
 \label{eq:perimeter}
\end{equation}
The corrected reference replaces the historical axis-recovery expression and normal-equation solver. Their exact difference, the historical outputs and source hashes are retained in the review archive. Historical human estimates are excluded from the results of the corrected method. The reference corrects section fitting and sensitivity recovery in isolation; the application pipeline has not been represented as fully patched or re-evaluated.

\subsection{Shoulder path and linear proxy}
At the engineering acromion level $Y_a=0.182H_m$, let $W,D$ be unexpanded traced width and profile depth. When $W,D>0$ and $D/W\leq0.6$, shoulder length is $S=\tfrac12P(W/2,D/2)$. This half-perimeter assumes that shoulder endpoints coincide with the transverse extrema of a symmetric section and that the posterior path is elliptical. Its contribution is a specified path model with explicit applicability checks; a tape trajectory can differ from these assumptions.

When this path cannot be computed, frontal silhouette tips are sought between $Y_s-20$\,mm and $Y_s+0.4(Y_e-Y_s)$, where $Y_e$ is the lower available elbow level. Missing elbows give an upper limit $Y_s+0.18(Y_h-Y_s)$. The trace retains the arm-fused prefix, applies a quadratic Savitzky--Golay window of at most nine rows, and selects each half-width maximum. A post-peak taper of at least 5\,mm on each side marks reliable tips. Their two-dimensional Euclidean separation gives projected breadth; a tip-to-landmark breadth ratio outside $[1,1.5]$ is advisory. Otherwise the fallback is the shoulder-landmark chord multiplied by $s_f$. The legacy fallback label can cover either breadth method, so the method must accompany the value. The output historically named inseam is $Y_{\rm ankle}-Y_{\rm hip}$, a vertical hip-to-ankle proxy without a corresponding tape reference in the historical development records.

\subsection{Display geometry and retained outputs}
Display rounding is applied to a separate volume. Each $XZ$ component is opened with a Euclidean distance-transform element of radius $0.9$ times its smaller half-extent, with a two-pixel minimum; an empty opening retains the component. Limb seeds satisfy frontal vertical thickness divided by profile-core depth $<0.5$, grow through ratios $<0.85$, and require frontal area at least $2000$\,mm$^2$. Limb $YZ$ openings are united with the rounded $XZ$ result. Distance-transform erosion uses $d_{\rm background}>r$ and dilation uses distance to the eroded set $\leq r$, with padding exceeding $r$. These elements differ from the rasterized OpenCV elements used for section measurement.

Junction repair attempts chessboard-distance closings of 20, 40, 60, and 100\,mm, intersecting the result with the unrounded profile volume. Processing stops at one component or residual components totalling at most $0.5\%$. The largest component is exported, with an advisory if more than $1\%$ is removed. Zero padding of $\lceil3\sigma_g\rceil+1$ voxels precedes Gaussian smoothing with $\sigma_g=1$ voxel and marching cubes at isovalue $0.5$ with spacing $h$~\cite{lorensen1987marching}. Origin and padding offsets are undone, and $(x,y,z)$ maps to $(x,z,y_{\max}-y)$ for a floor-based, vertical-$Z$ display frame. The largest mesh component is retained; a watertight mesh with negative signed volume is inverted. Twelve Taubin passes~\cite{taubin1995signal} use $\lambda=0.5$ and $\nu=-0.53$ under the applicable library convention. Optional simplification targets 150\,000 faces.

Saved outputs include normalized masks, conventional and display occupancy, and measurement records containing levels, fitted axes, rotations, and shoulder method. Advisory messages are largely logged and are not all serialized into the historical JSON. Reproducibility therefore requires retaining the diagnostic log with image/model identifiers and numerical settings. Dense Boolean occupancy consumes approximately $N_XN_YN_Z$ bytes; each full double-precision coordinate array adds eight times this storage. Chunked evaluation limits temporary arrays while leaving dense output storage unchanged.

The supplied research record contract separates endpoint identity from method identity. A posterior shoulder path cannot be relabelled as a projected breadth when a fallback occurs. Records carry attempt identifier, units, source branch, level, fallback reason, interpolation/depth-cap/waist flags, rejection status and a sensitivity availability reason. A rejected attempt has a null measurement and remains in the attempted-input denominator. Schema tests reject ambiguous legacy shoulder labels and unavailable sensitivity encoded as zero. This contract is a checked research interface; the unchanged application has not been claimed to emit it automatically.

\begin{table}[htbp]
\centering
\caption{Source geometry and defined outputs. Exported display surfaces supply no measurement in this table.}
\label{tab:branches}
\small
\begin{tabular}{p{0.34\linewidth}p{0.57\linewidth}}
\toprule
Output & Source and interpretation\\
\midrule
Chest-level girth & Profile-column occupancy; conventional hull only as a flagged fallback.\\
Waist-proxy/hip-landmark girth & Waist-adjusted conventional occupancy at the specified levels.\\
Posterior shoulder path & Half-ellipse model from frontal width and profile depth.\\
Projected shoulder breadth & Silhouette tips or landmark chord, identified separately from path.\\
Hip-to-ankle vertical proxy & Frontal landmark heights; no inseam interpretation.\\
\bottomrule
\end{tabular}
\end{table}

For the stated display coordinate map, the linear part has determinant $+1$ and preserves Euclidean distance; a known-object contract test verifies the mapping and floor offset. Mesh smoothing and junction repair alter the displayed surface, so this coordinate check supplies no anatomical surface-accuracy evidence. There is no full-resolution human runtime or peak-memory benchmark in this study; computational results are restricted to the synthetic kernels and grids actually executed.

\section{Conditional Sensitivity}
Let $g_b(q)$ denote the complete measurement for a fixed source branch $b$, where $q$ contains its mask coordinates, stature, levels and settings. First-order propagation, $\operatorname{Var}(g_b)\simeq J_b\Sigma_qJ_b^T$, requires a differentiable branch and a justified input covariance~\cite{jcgm2008}. The present verification does not estimate that covariance from human images. Component selection, rasterized morphology and rejection can make $g_b$ nonsmooth. The historical five-term root-sum-square index is excluded from the proposed reporting method; no confidence interval or coverage is claimed.

Within the ellipse model, write $q_e=(3a+b)(a+3b)$. Exact derivatives of the approximation in Eq.~\eqref{eq:perimeter} are
\begin{equation}
 P_a=\pi\left(3-\frac{3a+5b}{\sqrt{q_e}}\right),\qquad
 P_b=\pi\left(3-\frac{5a+3b}{\sqrt{q_e}}\right).
 \label{eq:derivative}
\end{equation}
For axis perturbations $(\Delta a,\Delta b)$, $\Delta P\simeq P_a\Delta a+P_b\Delta b$. This is a conditional ellipse calculation: a depth perturbation in a rotated section need not change only $b$. For the explicitly axis-aligned shoulder model $S=P(W/2,D/2)/2$, the corresponding derivatives are $\partial S/\partial W=P_a/4$ and $\partial S/\partial D=P_b/4$. A common stature perturbation gives $\Delta P/P\simeq\Delta H/H$ in the continuous scale-equivariant model; it is one shared contribution, and is not duplicated as independent view noise. A smoothly varying section also gives $\Delta C\simeq C'(Y)\Delta Y$ under a small level displacement. None of these identities supplies missing camera or path-model error.

The reference offers a conditional algebraic diagnostic. With normalized design $M$, residual $r=M\beta+\mathbf1$ and $n>5$, it computes $\hat\sigma^2=r^Tr/(n-5)$ and $\Sigma_\beta=\hat\sigma^2(M^TM)^{-1}$ through the SVD factors. This covariance assumes a fixed design and independent equal-variance algebraic residuals; it is not a statistical model of noisy contour coordinates. If $T=\operatorname{diag}(\rho^{-2},\rho^{-2},\rho^{-2},\rho^{-1},\rho^{-1})$, then $\Sigma_\phi=T\Sigma_\beta T^T$ and a numerical conic-perimeter Jacobian gives the diagnostic $\sqrt{J_\phi\Sigma_\phi J_\phi^T}$. Central differences use steps $\max(10^{-5}|\phi_i|,10^{-12})$ in the coefficient's own units and the corrected recovery in Eq.~\eqref{eq:recover}. Five-point fits, invalid perturbations and nonfinite covariance produce an explicit unavailable status with a reason and a null value. A successfully evaluated zero residual is a different state. No combined statistic is formed from unavailable terms.

The former curvature-to-area uncertainty and frontal-only shoulder budget are omitted: they were not derivatives of the actual outputs. The waist kernel is rerun with changed displacement, and its deleted voxel counts and centre-slice area changes are recorded. Section-fitting and perimeter responses to this deletion have not been evaluated. For a future complete branch sensitivity study, paired input perturbations must retain the same endpoint definition and separately record branch changes and rejected attempts. Empirical capture variation is required before such calculations can be interpreted as measurement uncertainty.

\section{Numerical Verification and Evidence Scope}
\subsection{Ellipse recovery and perimeter}
The reference test suite is deterministic, with random seed 20261021. An ellipse of semi-axes $(160,100)$\,mm is sampled at eleven nonuniform angles. Twenty-seven noiseless cases combine coordinate scales $0.01$, $1$ and $100$, rotations $0$, $0.47$ and $1.2$ radians, and translations $(0,0)$, $(20,-10)$ and $(10^6,-2\times10^6)$\,mm. The maximum relative semi-axis error is $3.35\times10^{-11}$. Additional exact aspect ratios 10 and 100, with full nonuniform and short-arc sampling, have relative axis errors below $4.92\times10^{-12}$; the largest scaled-design condition number is $1.58\times10^4$. These are floating-point checks on exact conics, rather than robustness guarantees for noisy partial contours. Fewer than five samples, repeated points, collinearity, nonfinite coordinates and a hyperbolic quadratic form exercise rejection.

For independent Gaussian coordinate noise of standard deviation $0.2$, $1$ or $3$\,mm, thirty trials per level use 180 nonuniform ellipse samples with fixed rotation $0.47$ radians. Perimeter RMSE against the exact generating ellipse is respectively $0.109$, $0.597$ and $1.759$\,mm. This experiment fixes the ellipse shape and noise family; it does not calibrate physical uncertainty or test silhouette segmentation. A five-point case and a case with invalid coefficient perturbations return unavailable sensitivity; ordinary exact fits retain an available, numerically near-zero algebraic diagnostic.

The perimeter reference is $4aE(m)$, where $m=1-b^2/a^2$ and $E$ is the complete elliptic integral in the parameter convention. Across 901 equally spaced ratios $b/a\in[0.1,1]$, the maximum relative error of Eq.~\eqref{eq:perimeter} is $8.42\times10^{-4}$, occurring at $b/a=0.1$. The analytic derivatives in Eq.~\eqref{eq:derivative} agree with central differences to less than $10^{-7}$ for tested axes $(100,100)$, $(160,100)$ and $(200,20)$\,mm. The approximation error and the fitting error are distinct; neither bounds errors caused by a non-elliptical body section.

\subsection{Profile-constraint experiments}
Controlled experiments use a $151\times200\times121$ grid with 4\,mm spacing over an 800\,mm-high box torso and prescribed limbs. The torso half-width is 100\,mm and its baseline depth is 120\,mm; horizontal limbs have 40\,mm vertical and sagittal extent. Binary frontal and profile silhouettes are exact projections of these primitives. Additional cases introduce a separate profile run, merged contamination, true depth variation, or displaced/vertical limbs. The oblique case changes the frontal limb centre by 0.6\,mm per lateral millimetre while keeping its prescribed vertical and sagittal thicknesses. Frozen historical profile-filtering and run-length kernels are used. Perception, torso-level selection, section morphology and girth fitting are excluded from this experiment.

\begin{table}[htbp]
\centering
\caption{Synthetic occupancy intersection-over-union against prescribed truth. H: conventional all-run hull; W: widest profile run without filtering; F: filtered interval without column cap; C: filtered interval with column cap. F and C use the numerical membership tolerance specified in the text.}
\label{tab:synthetic}
\small
\begin{tabular}{lrrrr}
\toprule
Case & H & W & F & C\\
\midrule
Horizontal limbs &0.94769&0.94769&0.94769&1.00000\\
Extra profile run &0.90682&0.94769&0.94769&1.00000\\
Merged artifact &0.81423&0.81423&0.94769&1.00000\\
True depth bulge &0.87533&0.87533&0.82112&0.86010\\
Sagittally displaced limb &0.90682&0.89322&0.89322&0.94107\\
Oblique constructed prism &0.94769&0.94769&0.94769&1.00000\\
Vertical limb &0.90914&0.90914&0.90914&0.90914\\
Smooth depth variation &0.96696&0.96696&0.81606&0.85810\\
\bottomrule
\end{tabular}
\end{table}

Table~\ref{tab:synthetic} separates favorable mechanisms from counterexamples. Selecting one run removes separated contamination; filtering removes the constructed merged artifact, and the column cap removes replicated limb depth. Genuine depth variation is also smoothed: on the smooth-depth case, full coupling reduces IoU from 0.96696 to 0.85810. A vertical limb receives no improvement, and a shared sagittal midpoint leaves displaced limb voxels incorrectly positioned. Perfect overlap in the oblique constructed prism follows its chosen vertical/depth thicknesses and does not establish robustness for anatomically oblique arms. Repeating the smooth-depth primitive with the same pixel-window settings at 2, 4 and 8\,mm spacing gives coupled IoUs 0.97738, 0.85810 and 0.80972. This variation demonstrates the dependence on physical filter bandwidth and sampling.

A numerical regression exposed floating-point drift when Savitzky--Golay filtering nominally constant, grid-aligned endpoints. Direct interval comparison removed boundary layers: the clean case's historical coupled IoU was 0.96862. The research kernel admits $z\in[z_- -\epsilon,z_++\epsilon]$ with $\epsilon=10^{-8}h$ in the same length units as $h$, giving the intended clean occupancy in Table~\ref{tab:synthetic}. This is an arithmetic tolerance, eight orders below one cell, and supplies no physical error allowance. Tests at $h=2$, 4 and 8\,mm recover boundaries drifted by $10^{-12}h$, preserve the loss caused by a genuine $10^{-5}h$ inset and add no adjacent voxel layer. Raw historical and tolerance-based results are both retained; this kernel revision is separate from the unchanged full application.

Waist tests use a prescribed half-width $w(Y)=80+0.001(Y-500)^2$\,mm and an otherwise box-like section. The rule gives $\kappa=0.002$\,mm$^{-1}$ and $d=20.8712$\,mm over a 400\,mm trace span. At zero displacement no voxels change. At half, full and capped stronger displacements, no voxels are added and respectively 4736, 8298 and 11984 are removed. Shifting the sagittal coordinates by 40\,mm changes the full-displacement deletion to 9060 voxels. The same silhouettes admit an unadjusted truth or one constructed from the deletion prior: the rule worsens overlap with the former and improves it with the latter. This counterexample shows why deletion bounds alone cannot establish anatomical benefit. A matched human waist/morphology ablation remains necessary.

The component-selection test also loses 729 voxels representing a deliberately disconnected valid limb. An empty profile is rejected. Missing pose and clipping after alignment have not been exercised end-to-end. The supplied tests account for their constructed cases; they do not estimate the attempted-capture failure rate of a human system. Analytic camera examples likewise illustrate a conditional effect: a $(160,100)$\,mm elliptical section viewed at $10^\circ$ yaw has an apparent frontal half-width 0.923\% smaller under orthographic projection. This is not a validated pose tolerance.

\subsection{Interpretation and requirements for human evaluation}
Historical development observations are preserved in the review archive and excluded from quantitative validation here. They precede the corrected conic reference and have undocumented reference endpoints and acquisition metadata. Source comments show that section rounding was checked against the development photographs; the available subjects therefore cannot be treated as an untouched test set. The present numerical settings are frozen for the supplied synthetic experiments. The package does not establish a corrected, integrated human pipeline or a held-out human result.

Physical validation requires tape readings at the same prescribed chest, waist-proxy and hip-landmark levels as the image outputs. Shoulder posterior-path endpoints must coincide with the model's transverse extrema; projected breadth must be evaluated separately. Posture, clothing, tape tension, operator, repeated readings and acquisition settings must be recorded. Reference intervals must be retained as intervals, and missing endpoints remain missing. Repeated sites within one participant do not increase the independent participant count. Agreement analysis should respect endpoint definitions and the sampling design~\cite{bland1986}. An independently acquired, ethically governed sample is needed to assess the utility of the profile and waist rules against established approaches, with all attempted captures retained when reporting failure rates.

\section{Conclusion}
The proposed pipeline couples a filtered profile interval to frontal column thickness and defines branch-specific geometric measurements. Its corrected conic reference passes reproducible synthetic recovery and derivative checks within the tested domains. The waist deletion and shoulder path are explicit shape hypotheses whose anatomical validity requires matched physical references. Numerical verification supports the stated computations; independent evidence is still required for an anthropometric accuracy claim or a demonstrated advantage over prior methods.

\paragraph{Availability and declarations.}
The accompanying review package contains the standalone manuscript, isolated corrected conic code, deterministic synthetic tests, results and an audit manifest. Historical application source, the corrected reference and synthetic kernels are distinguished in that manifest; a frozen, integrated public pipeline release is not claimed. Human images are outside this verification release. Author names, affiliations, contact details, funding and competing-interest declarations have not been supplied for this internal draft. The historical observations' consent, privacy arrangements and applicable ethics approval, exemption or institutional review status have not been established in the supplied material. These author-dependent requirements must be completed before journal submission; no clearance or permission is inferred.

\begin{thebibliography}{99}
\small
\setlength{\itemsep}{0.1em}
\setlength{\parskip}{0pt}
\bibitem{iso8559}
International Organization for Standardization, \emph{Size designation of clothes---Part 1: Anthropometric definitions for body measurement}, ISO 8559-1:2017, 2017. [Online]. Available: \url{https://www.iso.org/standard/61686.html}.

\bibitem{meunier2000}
P.~Meunier and S.~Yin, ``Performance of a 2D image-based anthropometric measurement and clothing sizing system,'' \emph{Applied Ergonomics}, vol.~31, no.~5, pp.~445--451, 2000. doi: \href{https://doi.org/10.1016/S0003-6870(00)00023-5}{10.1016/S0003-6870(00)00023-5}.

\bibitem{hung2004}
P.~C.-Y.~Hung, C.~P.~Witana, and R.~S.~Goonetilleke, ``Anthropometric measurements from photographic images,'' in \emph{Work with Computing Systems 2004}, H.~M.~Khalid, M.~G.~Helander, and A.~W.~Yeo, Eds. Kuala Lumpur: Damai Sciences, 2004, pp.~764--769. [Online]. Available: \url{https://www.ieda.ust.hk/dfaculty/ravi/papers/wwcs2004_1.pdf}.

\bibitem{aslam2017}
M.~Aslam, F.~Rajbdad, S.~Khattak, and S.~Azmat, ``Automatic measurement of anthropometric dimensions using frontal and lateral silhouettes,'' \emph{IET Computer Vision}, vol.~11, no.~6, pp.~434--447, 2017. doi: \href{https://doi.org/10.1049/iet-cvi.2016.0406}{10.1049/iet-cvi.2016.0406}.

\bibitem{laurentini1994visual}
A.~Laurentini, ``The visual hull concept for silhouette-based image understanding,'' \emph{IEEE Transactions on Pattern Analysis and Machine Intelligence}, vol.~16, no.~2, pp.~150--162, 1994.

\bibitem{loper2015smpl}
M.~Loper, N.~Mahmood, J.~Romero, G.~Pons-Moll, and M.~J.~Black, ``SMPL: A skinned multi-person linear model,'' \emph{ACM Transactions on Graphics}, vol.~34, no.~6, article~248, 2015.

\bibitem{varol2018bodynet}
G.~Varol, D.~Ceylan, B.~Russell, J.~Yang, E.~Yumer, I.~Laptev, and C.~Schmid, ``BodyNet: Volumetric inference of 3D human body shapes,'' in \emph{Proceedings of ECCV}, pp.~20--36, 2018.

\bibitem{ji2019}
Z.~Ji, X.~Qi, Y.~Wang, G.~Xu, P.~Du, X.~Wu, and Q.~Wu, ``Human body shape reconstruction from binary silhouette images,'' \emph{Computer Aided Geometric Design}, vol.~71, pp.~231--243, 2019. doi: \href{https://doi.org/10.1016/j.cagd.2019.04.019}{10.1016/j.cagd.2019.04.019}.

\bibitem{chen2025compact}
Q.~Chen, L.~Song, X.~Zhang, G.~Zhang, H.~Lu, J.~Liu, Y.~Pan, and W.~Luo, ``Accurate 3D anthropometric measurement using compact multi-view imaging,'' \emph{Measurement}, vol.~247, article~116777, 2025. doi: \href{https://doi.org/10.1016/j.measurement.2025.116777}{10.1016/j.measurement.2025.116777}.

\bibitem{lugaresi2019mediapipe}
C.~Lugaresi et al., ``MediaPipe: A framework for building perception pipelines,'' arXiv:1906.08172, 2019.

\bibitem{mediapipeguide}
Google AI Edge, ``Pose landmark detection guide.'' [Online]. Available: \url{https://ai.google.dev/edge/mediapipe/solutions/vision/pose_landmarker}. Accessed October 17, 2026.

\bibitem{rother2004grabcut}
C.~Rother, V.~Kolmogorov, and A.~Blake, ``GrabCut: Interactive foreground extraction using iterated graph cuts,'' \emph{ACM Transactions on Graphics}, vol.~23, no.~3, pp.~309--314, 2004.

\bibitem{schonemann1966generalized}
P.~H.~Sch\"onemann, ``A generalized solution of the orthogonal Procrustes problem,'' \emph{Psychometrika}, vol.~31, pp.~1--10, 1966.

\bibitem{savitzky1964smoothing}
A.~Savitzky and M.~J.~E.~Golay, ``Smoothing and differentiation of data by simplified least squares procedures,'' \emph{Analytical Chemistry}, vol.~36, no.~8, pp.~1627--1639, 1964.

\bibitem{graham1972efficient}
R.~L.~Graham, ``An efficient algorithm for determining the convex hull of a finite planar set,'' \emph{Information Processing Letters}, vol.~1, no.~4, pp.~132--133, 1972.

\bibitem{fitzgibbon1999direct}
A.~Fitzgibbon, M.~Pilu, and R.~B.~Fisher, ``Direct least square fitting of ellipses,'' \emph{IEEE Transactions on Pattern Analysis and Machine Intelligence}, vol.~21, no.~5, pp.~476--480, 1999.

\bibitem{ramanujan1914modular}
S.~Ramanujan, ``Modular equations and approximations to $\pi$,'' \emph{Quarterly Journal of Mathematics}, vol.~45, pp.~350--372, 1914.

\bibitem{lorensen1987marching}
W.~E.~Lorensen and H.~E.~Cline, ``Marching cubes: A high resolution 3D surface construction algorithm,'' \emph{ACM SIGGRAPH Computer Graphics}, vol.~21, no.~4, pp.~163--169, 1987.

\bibitem{taubin1995signal}
G.~Taubin, ``A signal processing approach to fair surface design,'' in \emph{Proceedings of SIGGRAPH}, pp.~351--358, 1995.

\bibitem{jcgm2008}
Joint Committee for Guides in Metrology, \emph{Evaluation of measurement data---Guide to the expression of uncertainty in measurement}, JCGM~100:2008. [Online]. Available: \url{https://doi.org/10.59161/JCGM100-2008E}.

\bibitem{bland1986}
J.~M.~Bland and D.~G.~Altman, ``Statistical methods for assessing agreement between two methods of clinical measurement,'' \emph{The Lancet}, vol.~327, no.~8476, pp.~307--310, 1986. doi: \href{https://doi.org/10.1016/S0140-6736(86)90837-8}{10.1016/S0140-6736(86)90837-8}.
\end{thebibliography}
\end{document}
